\section{基于持续同调特征的变点检验}\label{sec:method}

图~\ref{fig:flow}给出整体流程与统计冲突的来源：窗口是随机化单位，但重叠使原始观测被多个随机化单位共享。后续小节依次定义窗口构造、两种表示、最大扫描统计量与两种校准方法。

\begin{figure}[tbp]
\centering
\begin{subfigure}[t]{0.92\linewidth}
\centering
\caption{统计推断主流程}\label{fig:flow-a}
% 图 1(a)：重叠窗口拓扑扫描与校准的统计推断主流程（横向）。
% 依赖 main.tex 中定义的 flowblue/flowfill/m0fill/m1fill 配色与 tikz 库。
\begin{tikzpicture}[
  font=\footnotesize,
  box/.style={draw=flowblue, fill=flowfill, rounded corners=2pt,
              align=center, inner xsep=6pt, inner ysep=5pt, line width=0.6pt},
  arr/.style={-{Stealth[length=2.4mm]}, flowblue, line width=0.6pt},
]
\node[box] (raw) {原始序列\\ $X_1,\ldots,X_T$};
\node[box, right=4mm of raw] (win) {重叠滑动窗口\\ $I_1,\ldots,I_M$};
\node[box, right=4mm of win] (pc) {窗口点云\\ $P_i$};
\node[box, fill=m0fill, above right=3mm and 4mm of pc] (bpi) {持续同调 $\to$ 持续图\\ $\to$ BPI（800 维）};
\node[box, fill=m1fill, below right=3mm and 4mm of pc] (mom) {MOMENTS\\ （5/9 维）};
\coordinate (mid1) at ($(bpi.south east)!0.5!(mom.north east)$);
\node[box, right=4mm of mid1] (scan) {最大扫描\\ $T_M,\ \hat k$};

\node[box, fill=m0fill, below=15mm of scan, xshift=-19mm] (m0) {M0：全局随机置换重采样\\ $T_M^{*(1\ldots B)}$};
\node[box, fill=m1fill, below=15mm of scan, xshift=19mm] (m1) {M1：相关乘子重采样\\ （近似局部依赖）};

\coordinate (mid2) at ($(m0.south)!0.5!(m1.south)$);
\node[box, below=6mm of mid2] (ref) {参考分布 $\{T_M^*\}$\\ ＋ 观测 $T_M$};
\node[box, left=6mm of ref] (pv) {Monte Carlo\\ $p$ 值};
\node[box, left=6mm of pv] (dec) {拒绝 / 不拒绝\\ ＋ 位置估计 $\hat k$};

\draw[arr] (raw) -- (win);
\draw[arr] (win) -- (pc);
\draw[arr] (pc) -- (bpi);
\draw[arr] (pc) -- (mom);
\draw[arr] (bpi) -- (scan);
\draw[arr] (mom) -- (scan);
\draw[arr] (scan) -- (m0);
\draw[arr] (scan) -- (m1);
\draw[arr] (m0) -- (ref);
\draw[arr] (m1) -- (ref);
\draw[arr] (ref) -- (pv);
\draw[arr] (pv) -- (dec);
\end{tikzpicture}

\end{subfigure}

\vspace{1.2em}

\begin{subfigure}[t]{0.85\linewidth}
\centering
\caption{窗口重叠与共享观测机制}\label{fig:flow-b}
% 图 1(b)：窗口重叠与共享观测 + 全局随机置换破坏邻接。
\begin{tikzpicture}[
  font=\small,
  obs/.style={circle, draw=flowblue, fill=white, inner sep=1.2pt, line width=0.5pt},
  shared/.style={circle, draw=sharedred, fill=sharedred!25, inner sep=1.7pt, line width=0.6pt},
]
% 时间轴与观测
\foreach \i in {1,...,8} { \coordinate (X\i) at (\i, 0); }
\draw[flowblue, line width=0.5pt] (0.5,0) -- (8.5,0);
\foreach \i in {1,7,8} { \node[obs] at (X\i) {}; }
\foreach \i in {2,3,4,5,6} { \node[shared] at (X\i) {}; }
\foreach \i in {1,...,8} { \node[below=1.5mm, font=\footnotesize] at (X\i) {$X_{\i}$}; }

% 四个窗口（各覆盖 4 个连续观测，相邻窗口共享 3 个观测；与下方全局随机置换的 4 个窗口一致）
\draw[winblue, line width=1.3pt] (1,1.9) -- (4,1.9);
\draw[winblue, line width=0.8pt] (1,1.68) -- (1,2.12);
\draw[winblue, line width=0.8pt] (4,1.68) -- (4,2.12);
\node[right=1mm, font=\footnotesize] at (4,1.9) {窗口1};

\draw[winblue, line width=1.3pt] (2,1.35) -- (5,1.35);
\draw[winblue, line width=0.8pt] (2,1.13) -- (2,1.57);
\draw[winblue, line width=0.8pt] (5,1.13) -- (5,1.57);
\node[right=1mm, font=\footnotesize] at (5,1.35) {窗口2};

\draw[winblue, line width=1.3pt] (3,0.8) -- (6,0.8);
\draw[winblue, line width=0.8pt] (3,0.58) -- (3,1.02);
\draw[winblue, line width=0.8pt] (6,0.58) -- (6,1.02);
\node[right=1mm, font=\footnotesize] at (6,0.8) {窗口3};

\draw[winblue, line width=1.3pt] (4,0.25) -- (7,0.25);
\draw[winblue, line width=0.8pt] (4,0.03) -- (4,0.47);
\draw[winblue, line width=0.8pt] (7,0.03) -- (7,0.47);
\node[right=1mm, font=\footnotesize] at (7,0.25) {窗口4};

% 共享标注（置于窗口条上方）
\node[sharedred, font=\footnotesize] at (4.5,2.6) {共享观测（相邻窗口重叠 $W-s$ 个观测）};

% 全局随机置换示意（方块组居中于时间轴中心 x≈4.5，避免视觉偏左）
\foreach \i/\n in {1/1,2/2,3/3,4/4} {
  \node[draw=flowblue, fill=flowfill, minimum size=7mm, font=\footnotesize] at ({2.25+0.9*\i}, -2.3) {\n};
}
\node[left=2mm, font=\footnotesize] at (3.15,-2.3) {原始窗口序};
\foreach \i/\n in {1/3,2/1,3/4,4/2} {
  \node[draw=flowblue, fill=flowfill, minimum size=7mm, font=\footnotesize] at ({2.25+0.9*\i}, -3.8) {\n};
}
\node[left=2mm, font=\footnotesize] at (3.15,-3.8) {随机置换后};
\draw[-{Stealth[length=2mm]}, sharedred, line width=0.6pt] (3.6,-2.7) -- (3.6,-3.4);
\node[sharedred, right=2mm, font=\footnotesize] at (3.6,-3.05) {窗口级邻接被破坏};
\end{tikzpicture}

\end{subfigure}

\caption{基于持续同调特征的变点检验与校准流程。(a) 原始序列经重叠窗口映射为窗口点云，再分别经持续同调—向量化分支得到 BPI、经均值—协方差分支得到矩表示 MOMENTS，进入最大扫描得到检验统计量 $T_M$ 与位置估计 $\hat k$；随后 M0 无约束窗口级随机置换重采样、M1 相关乘子重采样分别构造参考分布，并通过 Monte Carlo $p$ 值形成拒绝决策。(b) 重叠机制：相邻窗口共享原始观测（红色），无约束窗口级随机置换重排窗口顺序时原有邻接被破坏，这是后文交换性失效机制的直观来源。}\label{fig:flow}
\end{figure}

\subsection{窗口构造与局部点云}\label{sec:windows}

给定观测序列 $X_1,\ldots,X_T$，取窗口宽度 $W$ 和步长 $s$，并令
\begin{equation}\label{eq:windows}
 a_i=(i-1)s+1,\qquad
 I_i=\{a_i,\ldots,a_i+W-1\},\qquad
 M=\left\lfloor\frac{T-W}{s}\right\rfloor+1.
\end{equation}
本文考虑 $1\le s\le W$，相邻窗口的名义重叠率为 $1-s/W$。对于二维观测，窗口点云直接定义为 $P_i=\{X_t:t\in I_i\}$。对于标量观测，在每个窗口内部构造三维时延点云
\begin{equation}\label{eq:delay-cloud}
 P_i=\bigl\{(X_t,X_{t+3},X_{t+6}):t=a_i,\ldots,a_i+W-7\bigr\},
\end{equation}
因而时延坐标不会跨越窗口边界。点云在计算持续同调前按固定规则进行精确去重；所有坐标尺度、网格、核宽和截断参数均预先设定，且不利用整条序列共同估计的自适应参数。该局部性条件将在第~\ref{sec:theory} 节用于界定窗口特征的依赖范围。

\subsection{固定网格持续强度表示}\label{sec:bpi}

对每个点云 $P_i$，在系数域 $\mathbb F_2$ 上按欧氏距离计算 Vietoris--Rips 持续同调，并保留 0 维和 1 维同调通道中有限且寿命为正的 birth--death 对。过滤采用边长阈值的直径约定：单形 $\sigma$ 的过滤值为其顶点间最大欧氏距离，即 $\operatorname{diam}(\sigma)$；尺度 $\varepsilon$ 下纳入所有满足 $\operatorname{diam}(\sigma)\le\varepsilon$ 的单形，即 $\operatorname{VR}_\varepsilon(P)=\{\sigma\subseteq P:\operatorname{diam}(\sigma)\le\varepsilon\}$。过滤不设数据驱动的截断上限，持续图坐标直接以欧氏距离尺度报告。由于有限欧氏点云的 Vietoris--Rips 过滤在过滤值达到点云直径时形成完全复形，所有 1 维同调类在此时均已消亡，仅最后一个 0 维连通分量的 death 为 $\infty$；实现中只丢弃该显式无限 death 分量，其余有限类全部保留。相关计算框架和有限维表示可分别参见\cite{Bauer2021,Adams2017}。

记 $K_i(\varepsilon)=\operatorname{VR}_\varepsilon(P_i)$。随过滤尺度 $\varepsilon$ 增大，复形形成嵌套序列 $K_i(\varepsilon)\subseteq K_i(\varepsilon')$（$\varepsilon\le\varepsilon'$）。零维同调记录连通分支，一维同调记录未被二维单形填充的独立循环；贝蒂数 $\beta_q(K_i(\varepsilon))=\dim H_q(K_i(\varepsilon);\mathbb F_2)$ 表示该尺度下相应同调空间的维数。持续同调进一步记录这些类随过滤尺度的出现与消失，由持续图中的 $(b,d)$ 概括，其中 $b$ 与 $d$ 分别为出生尺度与消亡尺度，$d-b$ 是过滤尺度上的寿命，而非真实时间轴上的持续时长。记第 $q\in\{0,1\}$ 个同调通道的持续图为
\[
 \mathcal D_i^{(q)}=\{(b_{ij}^{(q)},d_{ij}^{(q)})\}_j.
\]
给定场景固定的尺度上界 $c$，将每个有限持续图点转换为归一化 birth--lifetime 坐标。尺度上界 $c$ 仅用于该坐标缩放与截断，不用于提前终止过滤；过滤本身以 $+\infty$ 为阈值完整计算。
\begin{equation}\label{eq:bpi-normalization}
 u_{ij}^{(q)}=\min\left(\frac{b_{ij}^{(q)}}{c},1\right),\qquad
 v_{ij}^{(q)}=\min\left(\frac{d_{ij}^{(q)}-b_{ij}^{(q)}}{c},1\right).
\end{equation}
令 $\Delta=1/20$、$g_r=(r+1/2)\Delta$（$r=0,\ldots,19$）以及 $\sigma=0.05$。第 $q$ 个通道在网格中心 $(g_r,g_\ell)$ 上的持续强度坐标定义为
\begin{equation}\label{eq:bpi}
 Z_{i,r\ell}^{(q)}
 =\frac{\Delta}{2\pi\sigma^2}
 \sum_j v_{ij}^{(q)}
 \exp\left\{-\frac{(g_r-u_{ij}^{(q)})^2+(g_\ell-v_{ij}^{(q)})^2}{2\sigma^2}\right\}.
\end{equation}
每个同调通道产生 $20\times20=400$ 个坐标。按 0 维、1 维同调通道的顺序向量化并连接，再除以 $\sqrt 2$（该缩放使两个同调通道对总平方距离的贡献取等权平均），得到 800 维特征。此处涉及三类维数：同调维数 $q\in\{0,1\}$、点云所在空间维数 $d\in\{2,3\}$ 以及向量表示维数 $p=800$；$800$ 维特征不意味着本文研究了 $800$ 维拓扑。
\begin{equation}\label{eq:bpi-vector}
 Y_i^{\mathrm{BPI}}
 =\frac{1}{\sqrt2}
 \begin{bmatrix}
   \operatorname{vec}(Z_i^{(0)})\\
   \operatorname{vec}(Z_i^{(1)})
 \end{bmatrix}.
\end{equation}
空通道记为 400 维零向量。N-AR 时延点云取 $c=6$，圆形模型及其备择取 $c=2.5$。本文将每个持续图中以寿命加权的高斯核表面在固定网格中心进行评价，并将两个同调通道的结果拼接为向量，沿用实现标识 BPI。该表示与持续图像共享“持续图--加权核表面--有限维表示”的思路，但其离散化采用网格中心采样，区别于逐像素积分\cite{Adams2017}。BPI 仅标识本文的固定实现，不另行宣称一种标准 Betti 不变量，也不构成对随机持续图总体强度密度的估计\cite{WuKimRinaldo2023}。

\subsection{矩表示}\label{sec:moments}

为构造低阶统计摘要作为对照，设窗口点云为 $P_i=\{x_{ij}\}_{j=1}^{n_i}\subset\mathbb R^d$，并定义
\begin{equation}\label{eq:moments}
 \bar x_i=\frac1{n_i}\sum_{j=1}^{n_i}x_{ij},\qquad
 S_i=\frac1{n_i}\sum_{j=1}^{n_i}(x_{ij}-\bar x_i)(x_{ij}-\bar x_i)^\top.
\end{equation}
MOMENTS 特征依次连接均值、协方差对角元以及按固定字典序排列的 $\sqrt2(S_i)_{ab}$（$a<b$）：
\begin{equation}\label{eq:moments-vector}
 Y_i^{\mathrm{MOM}}
 =\bigl(\bar x_i^\top,(S_i)_{11},\ldots,(S_i)_{dd},
 \{\sqrt2(S_i)_{ab}:a<b\}\bigr)^\top.
\end{equation}
其维数为 $d(d+3)/2$，故二维和三维点云分别得到 5 维和 9 维特征。$\sqrt2$ 权重保证任意两个窗口 $i,j$ 满足
\begin{equation}\label{eq:moments-frob}
 \|Y_i^{\mathrm{MOM}}-Y_j^{\mathrm{MOM}}\|_2^2
 =\|\bar x_i-\bar x_j\|_2^2+\|S_i-S_j\|_F^2.
\end{equation}
该表示主要概括一、二阶矩；有限窗口和序列依赖仍可能使其响应某些形状变化，因此本文不预设其对所有拓扑变化均不敏感。

\subsection{最大扫描统计量与位置估计}\label{sec:scan}

本文采用的最大扫描统计量属于变点分析中标准化部分和型统计量的常见构造\cite{CsorgoHorvath1997}。对任一窗口特征序列 $Y_1,\ldots,Y_M\in\mathbb R^p$，令
\begin{equation}\label{eq:bridge}
 U_M(k)=\frac1{\sqrt M}
 \left(\sum_{i=1}^kY_i-\frac{k}{M}\sum_{i=1}^M Y_i\right),
 \qquad t_k=\frac{k}{M},
\end{equation}
并定义标准化扫描曲线与最大统计量
\begin{equation}\label{eq:scan}
 Q_M(k)=\frac{\|U_M(k)\|_2^2}{t_k(1-t_k)},\qquad
 T_M=\max_{k\in\mathcal K_M}Q_M(k),
\end{equation}
其中
\[
 \mathcal K_M=\{\lceil0.15M\rceil,\ldots,M-\lceil0.15M\rceil\},
\]
比例 $0.15$ 为确认性分析前固定的边界保护比例，用于避免左右片段过短。
等价地，若 $\bar Y_{1:k}$ 和 $\bar Y_{k+1:M}$ 分别表示左右两段样本均值，则
\begin{equation}\label{eq:scan-equivalent}
 Q_M(k)=\frac{k(M-k)}{M}
 \|\bar Y_{1:k}-\bar Y_{k+1:M}\|_2^2.
\end{equation}
因此，该扫描统计量直接响应左右两段窗口特征的均值之差，检验对象是特征均值层面的变化；它不构成对任意原始分布或任意拓扑变化都具有一致检出能力的通用检验。

对 BPI 表示，令 $V_i^{(q)}=\operatorname{vec}(Z_i^{(q)})$ 为第 $q$ 个同调通道的未缩放向量，并分别对每个通道使用相同的桥接与标准化，得到 $Q_M^{(q)}(k)$。由 $Y_i^{\mathrm{BPI}}=\frac1{\sqrt2}[V_i^{(0)};V_i^{(1)}]$ 及平方欧氏范数的线性性，对每个候选位置 $k$ 有
\[
 Q_M^{\mathrm{BPI}}(k)=\frac12 Q_M^{(0)}(k)+\frac12 Q_M^{(1)}(k).
\]
因此总扫描同时受连通分支持续信息与循环持续信息影响。由于最大值的位置通常不同，$\max_k\tfrac12\bigl(Q_M^{(0)}(k)+Q_M^{(1)}(k)\bigr)$ 一般不等于两个通道最大值的平均 $\tfrac12\max_k Q_M^{(0)}(k)+\tfrac12\max_k Q_M^{(1)}(k)$。本文实验将两个同调通道合并报告，未单独分离零维与一维同调的贡献。

为使浮点并列情形具有确定性，令 $\epsilon_{64}=2^{-52}$，将满足
\[
 Q_M(k)\ge T_M-64\epsilon_{64}\max(1,|T_M|)
\]
的连续候选段定义为合格平台。位置估计选取最长平台；长度并列时取最左平台，再以该平台的下中位数作为 $\widehat k$。由于 $k$ 表示第 $k$ 与第 $k+1$ 个窗口之间的分割，映射到原始时间轴后的连续边界估计为
\begin{equation}\label{eq:location-map}
 \widehat\tau=\left(\widehat k-\frac12\right)s+\frac{W+1}{2}.
\end{equation}
该映射将第 $\widehat k$ 与第 $\widehat k+1$ 个窗口之间的分割点对应到二者中心之间的中点。若特征序列精确恒定，则规定 $T_M=0$。所有校准方法均使用相同的观测统计量、候选集合和位置规则。

\subsection{校准方法}\label{sec:calibration}

朴素置换方法记为 M0。M0 从排列群 $\mathfrak S_M$ 上的均匀分布中独立抽取随机排列；下文称为全局随机置换，以区别于块置换或局部置换。条件于观测到的窗口特征，M0 对每个随机排列序列重新计算式\eqref{eq:scan}。该方法保留经验边际分布，却不保留原始窗口顺序中的局部依赖。

依赖乘子自助法记为 M1。令 $\bar Y=M^{-1}\sum_{i=1}^M Y_i$、$Z_i=Y_i-\bar Y$，并取
\begin{equation}\label{eq:dmb-bandwidth}
 b=\max\left\{\left\lfloor\frac Ws\right\rfloor+1,
                 \left\lceil M^{1/3}\right\rceil\right\}.
\end{equation}
乘子新息 $\varepsilon_1,\ldots,\varepsilon_{M+b-1}$ 为独立标准正态，由与观测样本分离的随机流生成，因而与观测样本相互独立；不同内层重复使用各自的随机流独立生成。定义
\begin{equation}\label{eq:dmb-multiplier}
 \xi_i=b^{-1/2}\sum_{j=0}^{b-1}\varepsilon_{i+j},
 \qquad i=1,\ldots,M.
\end{equation}
由此
\begin{equation}\label{eq:dmb-cov}
 \operatorname{Cov}(\xi_i,\xi_{i+h})
 =\left(1-\frac{|h|}{b}\right)_+.
\end{equation}
M1 不作循环回绕，并对乘子加权序列 $\{\xi_iZ_i\}_{i=1}^M$ 应用与观测样本相同的桥接扫描。该构造属于依赖感知的参考近似；dependent wild bootstrap 和依赖乘子过程的一般理论可参见\cite{Shao2010,BucherKojadinovic2016}，其结论不直接构成本文完整 BPI 最大扫描的有限样本有效性定理。

两种方法均使用 $B=999$ 次重采样。若
\[
 K=\sum_{r=1}^{B}\mathbf 1\{T_{M,r}^*\ge T_M\},
\]
则 Monte Carlo $p$ 值定义为
\begin{equation}\label{eq:mc-pvalue}
 p=\frac{K+1}{B+1}=\frac{K+1}{1000}.
\end{equation}
本文以 $p\le0.05$ 拒绝零假设，等价于 $K\le49$。加一形式避免有限重采样下报告零 $p$ 值\cite{PhipsonSmyth2010}，但不恢复重叠窗口的交换性，也不替代对 M1 参考近似有效性的评估；随机排列检验的有限样本有效性仍需相应的群不变性条件\cite{HemerikGoeman2018}。
